\section{Билет 20. Вывод одномерного волнового дифференциального уравнения для плоских волн.}

\begin{definition}
\textbf{Одномерным волновым уравнением} называется линейное дифференциальное уравнение в частных производных второго порядка, описывающее распространение возмущений вдоль одной пространственной координаты:
\begin{equation}
    \frac{\partial^2 \xi}{\partial t^2} = v^2 \frac{\partial^2 \xi}{\partial x^2}. \label{eq:wave_eq_1d}
\end{equation}
Здесь $\xi(x, t)$ --- смещение частиц среды, $v$ --- фазовая скорость распространения волны.
\end{definition}

\begin{theorem}[Вывод волнового уравнения]
Любая плоская волна, распространяющаяся вдоль оси $x$ без дисперсии и поглощения, может быть представлена в виде суперпозиции двух бегущих волн: одной в положительном направлении, другой в отрицательном. Рассмотрим произвольную функцию, описывающую волну, бегущую в сторону возрастания $x$:
\begin{equation}
    \xi(x, t) = f(x - vt), \label{eq:dalambert}
\end{equation}
где $f$ --- дважды дифференцируемая функция произвольного вида (не обязательно гармоническая). Введём вспомогательную переменную $u = x - vt$. Тогда $\xi = f(u)$.

Вычислим частные производные по времени, используя цепное правило:
\begin{equation}
    \frac{\partial \xi}{\partial t} = \frac{df}{du} \frac{\partial u}{\partial t} = -v f'(u), \quad
    \frac{\partial^2 \xi}{\partial t^2} = (-v) \frac{d}{du}[f'(u)] \frac{\partial u}{\partial t} = v^2 f''(u). \label{eq:deriv_t}
\end{equation}
Теперь вычислим частные производные по координате:
\begin{equation}
    \frac{\partial \xi}{\partial x} = \frac{df}{du} \frac{\partial u}{\partial x} = f'(u), \quad
    \frac{\partial^2 \xi}{\partial x^2} = \frac{d}{du}[f'(u)] \frac{\partial u}{\partial x} = f''(u). \label{eq:deriv_x}
\end{equation}
Сравнивая \eqref{eq:deriv_t} и \eqref{eq:deriv_x}, видим, что они связаны соотношением:
\begin{equation}
    \frac{\partial^2 \xi}{\partial t^2} = v^2 \frac{\partial^2 \xi}{\partial x^2}.
\end{equation}
Аналогичное соотношение получается для волны $\xi = g(x + vt)$, бегущей в противоположном направлении. В силу линейности уравнения, сумма решений также является решением. Таким образом, любая плоская волна удовлетворяет уравнению \eqref{eq:wave_eq_1d}.
\end{theorem}

\begin{remark}
Уравнение \eqref{eq:wave_eq_1d} является фундаментальным для волновой физики. Его важность заключается в следующем:
\begin{enumerate}
    \item Если при анализе любой физической системы (упругая среда, электромагнитное поле, квантовые вероятности) мы приходим к уравнению вида \eqref{eq:wave_eq_1d}, мы можем сразу утверждать, что в системе существуют волновые процессы со скоростью $v$.
    \item Уравнение линейно, поэтому для него справедлив \textbf{принцип суперпозиции}: если $\xi_1$ и $\xi_2$ --- решения, то их линейная комбинация $C_1\xi_1 + C_2\xi_2$ также является решением.
    \item Общее решение уравнения (формула Даламбера) имеет вид $\xi(x,t) = f(x-vt) + g(x+vt)$, что физически означает независимое распространение волн в противоположных направлениях без искажения формы (в отсутствие дисперсии и затухания).
\end{enumerate}
\end{remark}
